\scanpage{16}
% Foliotation archivistique au crayon, en bas à gauche : 15.
Passant à la limite projective on trouve
\[
\mathcal G_\infty=\operatorname{Aut}_k(\mathcal X_\infty\to\mathcal X)\simeq\Gamma\cdot T_\infty(\overline{\mathbb Q}^{*})
\]
\inserted{$\mathcal X_\infty$ : rev. abélien étale maximal de $\mathcal X$}

\ldots

C) $\operatorname{Card}S=3$\qquad
$\mathbb G\xrightarrow{\sim}\unsure{\mathfrak S_S|\mathbb G_S}\ (\simeq\unsure{\mathfrak S_3|\mathbb G_3})$

Il y a une façon canonique de réduire le corps de base à $k$, de façon que $S$ devienne rat$/k$. On est conduit ici à la tour de Fermat :
\[
\mathcal G_n=\bigl(\Gamma\times
\underbrace{\unsure{\mathfrak S_S|\mathbb G_S}\bigr)\cdot\pi_n}_{\mathbb G_{(n)}}
\qquad \pi_n=\mu_n^S/\mu_n
\]
\deleted{$\underbrace{\phantom{\Gamma\times\mathfrak S_S}}_{\unsure{\mathbb G_{(n)}|\mathcal G_{(n)}}}$}
% L'accolade biffée est sous le premier facteur ; la nouvelle accolade commence sous le second facteur et finit sous pi_n.

\ldots

D) $\operatorname{Card}S\geq4$. Soit $s=(s_0,s_1,s_\infty)$ \inserted{un triplet de} trois points distincts de $S$, on les utilise pour définir un iso
\[
\widehat{\mathcal X}\underset{\varphi_s}{\simeq}\mathbb P^1_{\overline{\mathbb Q}}
\]
transformant $s_0$ en $0$, $s_1$ en $1$, $s_\infty$ en $\infty$. Les autres points \deleted{sont} \inserted{de $S$ s'identifient} ainsi à des élts de $\overline{\mathbb Q}\setminus\{0,1\}$, soit $K_s\subset\overline{\mathbb Q}$
